%==============================================================================
%  FinPilot -- Personal Finance & Investment Dashboard
%  Full-Stack Development Internship -- Codveda Technologies
%  Internship Documentation Report
%
%  Compile with:  pdflatex  (run twice for the table of contents), or xelatex.
%
%    pdflatex FinPilot-Codveda-Internship-Documentation.tex
%    pdflatex FinPilot-Codveda-Internship-Documentation.tex
%
%  Images are referenced from the project's ../screenshots/ folders.
%  See report/README.md for the full image mapping.
%==============================================================================
\documentclass[11pt,a4paper]{report}

%------------------------------------------------------------------ packages ---
\usepackage[T1]{fontenc}
\usepackage[utf8]{inputenc}
\usepackage{lmodern}
\usepackage[a4paper,margin=1in,headheight=15pt]{geometry}
\usepackage{graphicx}
\usepackage{xcolor}
\usepackage{titlesec}
\usepackage{fancyhdr}
\usepackage{booktabs}
\usepackage{array}
\usepackage{longtable}
\usepackage{tabularx}
\usepackage{enumitem}
\usepackage{caption}
\usepackage{listings}
\usepackage{float}
\usepackage[hidelinks]{hyperref}

%-------------------------------------------------------------------- colors ---
\definecolor{brandprimary}{HTML}{1F3A5F}   % deep navy
\definecolor{brandaccent}{HTML}{2E86AB}    % teal blue
\definecolor{brandlight}{HTML}{EAF2F8}     % light panel
\definecolor{codebg}{HTML}{F5F7FA}
\definecolor{codekw}{HTML}{1F3A5F}
\definecolor{codecomment}{HTML}{6A737D}
\definecolor{codestring}{HTML}{2E7D32}
\definecolor{rulegray}{HTML}{B0BEC5}

%---------------------------------------------------------------- hyperlinks ---
\hypersetup{
  colorlinks=true,
  linkcolor=brandprimary,
  urlcolor=brandaccent,
  citecolor=brandaccent,
  pdftitle={FinPilot -- Codveda Full-Stack Internship Documentation},
  pdfauthor={Sushil Bhatta}
}

%----------------------------------------------------------- section styling ---
\titleformat{\chapter}[display]
  {\normalfont\bfseries\color{brandprimary}}
  {\filright\Large\color{brandaccent}\MakeUppercase{\chaptertitlename}\ \thechapter}
  {6pt}
  {\Huge}
\titlespacing*{\chapter}{0pt}{10pt}{24pt}

\titleformat{\section}
  {\normalfont\Large\bfseries\color{brandprimary}}{\thesection}{0.6em}{}
\titleformat{\subsection}
  {\normalfont\large\bfseries\color{brandaccent}}{\thesubsection}{0.6em}{}

%--------------------------------------------------------- header and footer ---
\pagestyle{fancy}
\fancyhf{}
\renewcommand{\headrulewidth}{0.4pt}
\renewcommand{\footrulewidth}{0.4pt}
\fancyhead[L]{\small\color{brandprimary}\textbf{FinPilot}}
\fancyhead[R]{\small\color{brandprimary}Codveda Full-Stack Internship}
\fancyfoot[L]{\small Sushil Bhatta}
\fancyfoot[C]{\small FinPilot -- Personal Finance \& Investment Dashboard}
\fancyfoot[R]{\small\thepage}

\fancypagestyle{plain}{%
  \fancyhf{}
  \renewcommand{\headrulewidth}{0pt}
  \renewcommand{\footrulewidth}{0.4pt}
  \fancyfoot[C]{\small\thepage}
}

%------------------------------------------------------------ listing styling ---
\lstdefinestyle{finpilot}{
  backgroundcolor=\color{codebg},
  basicstyle=\ttfamily\footnotesize,
  keywordstyle=\color{codekw}\bfseries,
  commentstyle=\color{codecomment}\itshape,
  stringstyle=\color{codestring},
  numberstyle=\tiny\color{codecomment},
  numbers=left,
  numbersep=8pt,
  frame=single,
  rulecolor=\color{rulegray},
  breaklines=true,
  showstringspaces=false,
  tabsize=2,
  captionpos=b,
  xleftmargin=14pt,
  framexleftmargin=14pt
}
\lstset{style=finpilot}

%-------------------------------------------------------------- table helper ---
\newcolumntype{L}[1]{>{\raggedright\arraybackslash}p{#1}}
\renewcommand{\arraystretch}{1.25}

% Consistent figure helper
\newcommand{\finfig}[3]{%
  \begin{figure}[H]
    \centering
    \includegraphics[width=#2\linewidth,keepaspectratio]{#1}
    \caption{#3}
  \end{figure}%
}

\captionsetup{font=small,labelfont={bf,color=brandprimary}}

%==============================================================================
\begin{document}

%------------------------------------------------------------------ COVER ------
\begin{titlepage}
  \thispagestyle{empty}
  \begin{center}
    \vspace*{0.6cm}

    {\color{brandaccent}\rule{\linewidth}{2pt}}\\[0.4cm]
    {\Huge\bfseries\color{brandprimary} FinPilot}\\[0.35cm]
    {\LARGE\color{brandprimary} Personal Finance \& Investment Dashboard}\\[0.30cm]
    {\color{brandaccent}\rule{\linewidth}{2pt}}\\[1.4cm]

    {\Large\bfseries Full-Stack Development Internship}\\[0.5cm]
    {\large Internship Project Documentation}\\[1.6cm]

    \begin{tabular}{rl}
      \textbf{Organization:} & Codveda Technologies \\[0.35cm]
      \textbf{Project:}      & FinPilot -- Personal Finance \& Investment Dashboard \\[0.35cm]
      \textbf{Project Type:} & Full-Stack Web Application (MERN) \\[0.35cm]
      \textbf{Submitted by:} & Sushil Bhatta \\
    \end{tabular}

    \vfill

    \begin{tabular}{rl}
      \textbf{Live App:} & \url{https://finpilot-oiek-q4zxkmfqt-sushil-bhattas-projects.vercel.app/} \\[0.2cm]
      \textbf{API Base:} & \url{https://finplot-rhir.onrender.com} \\[0.2cm]
      \textbf{Repository:} & \url{https://github.com/Sushil-Bhatta-sb/finpilot} \\
    \end{tabular}

    \vspace{1.2cm}
    {\color{brandaccent}\rule{\linewidth}{1pt}}
  \end{center}
\end{titlepage}

%----------------------------------------------------------- TABLE OF CONTENTS -
\pagenumbering{roman}
\tableofcontents
\clearpage
\pagenumbering{arabic}

%==============================================================================
\chapter{Project Overview}
%==============================================================================

\section{Introduction}
FinPilot is a full-stack personal finance dashboard for tracking income,
expenses, budgets, savings goals, investments, and transactions. It is designed
for individual users who want a single, privacy-respecting dashboard to monitor
their personal finances. The application provides a consolidated view of income
and expenses, budget progress and alerts, savings goals, an investments tracker,
transaction history with export options, and role-based admin tools.

\section{Purpose}
The goal of FinPilot is to give an individual user one place to understand where
their money comes from and where it goes. Every financial record is scoped to the
authenticated user, so a user can only view and modify their own data. Admin
users additionally receive platform-level management and reporting tools.

\section{Major Features}
\begin{itemize}[leftmargin=1.4em]
  \item \textbf{Authentication} --- JWT-based auth with bcrypt password hashing
        and role-based access control.
  \item \textbf{Financial Tracking} --- Create and manage income, expenses,
        budgets (with exceed alerts), savings goals, and investments.
  \item \textbf{Insights} --- Interactive dashboard with charts and category
        breakdowns.
  \item \textbf{Real-Time} --- Socket.io powered per-user notifications and
        toast alerts.
  \item \textbf{Admin Tools} --- User management, suspend accounts, and platform
        reporting for admins.
  \item \textbf{Data Export} --- Export transaction history to CSV and PDF.
\end{itemize}

\section{Architecture}
FinPilot follows a classic three-tier full-stack architecture. The React
frontend communicates with the Express REST API over HTTP, and the API persists
data in MongoDB through Mongoose. A Socket.io channel runs alongside the REST API
to deliver real-time notifications to individual users.

\begin{center}
\vspace{0.3cm}
{\large\bfseries\color{brandprimary} Frontend \;$\longrightarrow$\; Backend API \;$\longrightarrow$\; Database}\\[0.2cm]
{\small React (Vercel) \quad$\longrightarrow$\quad Node/Express (Render) \quad$\longrightarrow$\quad MongoDB Atlas}
\vspace{0.3cm}
\end{center}

\section{Technology Stack}
\begin{table}[H]
\centering
\begin{tabularx}{\linewidth}{>{\bfseries}l X}
\toprule
Layer & Technologies \\
\midrule
Frontend   & React, TypeScript, Vite, React Router, Axios, Recharts, Socket.io Client \\
Backend    & Node.js, Express, JWT, bcryptjs, Socket.io, PDFKit \\
Database   & MongoDB (Mongoose) \\
Deployment & Vercel (frontend), Render (backend), MongoDB Atlas \\
\bottomrule
\end{tabularx}
\caption{FinPilot technology stack by layer.}
\end{table}

\section{Authentication Behavior}
The dashboard and general pages can initially be viewed without logging in ---
the dashboard statistics endpoint returns a guest-safe payload when the request
is unauthenticated. Authentication is required only when accessing personal or
database-specific functionality. Protected operations include all user-specific
financial data (income, expenses, budgets, savings, investments, transactions,
categories) and profile/account operations. These protections are enforced on
the backend as well, so the API rejects unauthenticated access to private routes
regardless of the client used.

\section{Project Structure}
\begin{lstlisting}[language=bash,caption={FinPilot repository structure.}]
finpilot/
|- backend/
|  |- server.js
|  |- render.yaml
|  \- src/
|     |- app.js
|     |- socket.js
|     |- config/
|     |- controller/
|     |- middleware/
|     |- models/
|     \- routes/
|- frontend/
|  |- package.json
|  \- src/
|     |- api/
|     |- components/
|     |- context/
|     |- pages/
|     \- types/
\- docs/
   |- level-1/
   |- level2/
   \- level-3/
\end{lstlisting}

%==============================================================================
\chapter{Level 1 --- Basic}
%==============================================================================

\section{Task 1 --- Development Environment Setup}

\subsection{Objective}
Set up and verify a working local development environment for building the
FinPilot full-stack application, and initialize version control.

\subsection{Tools Installed}
\begin{table}[H]
\centering
\begin{tabularx}{\linewidth}{>{\bfseries}l X}
\toprule
Tool & Version / Detail \\
\midrule
Node.js      & v20.x \\
npm          & v10.x \\
Git          & v2.x \\
Code Editor  & VS Code \\
Database     & MongoDB Atlas (cloud-hosted) \\
\bottomrule
\end{tabularx}
\caption{Development environment toolchain.}
\end{table}

\subsection{Version Control}
The project repository is hosted on GitHub at
\url{https://github.com/Sushil-Bhatta-sb/finpilot}. Basic Git commands were
practiced as part of the setup: \texttt{git init}, \texttt{git add .},
\texttt{git commit -m "message"}, and \texttt{git push origin main}.

\subsection{Verification}
\finfig{../screenshots/level-1/tool-versions.png}{0.8}{Verifying installed Node.js, npm, and Git versions.}
\finfig{../screenshots/level-1/mongo-connected.png}{0.8}{Successful connection to the MongoDB Atlas database.}

\subsection{Contribution to FinPilot}
This task established the toolchain, cloud database, and version control that
every later task builds upon, ensuring a reproducible environment for developing
and deploying the application.

\section{Task 2 --- REST API}

\subsection{Objective}
Build a simple REST API providing full CRUD operations for two core resources of
the finance domain: \textbf{Income} and \textbf{Expense}.

\subsection{Technologies Used}
Node.js, Express, and MongoDB via the Mongoose ODM.

\subsection{Implementation}
A REST API was implemented with full CRUD for the Income and Expense resources,
input validation, and centralized error handling. The endpoints follow REST
conventions and return JSON responses.

\begin{table}[H]
\centering
\begin{tabularx}{\linewidth}{l l X}
\toprule
\textbf{Method} & \textbf{Endpoint} & \textbf{Description} \\
\midrule
GET    & /api/income      & Get all income \\
POST   & /api/income      & Create income entry \\
PUT    & /api/income/:id  & Update income entry \\
DELETE & /api/income/:id  & Delete income entry \\
\midrule
GET    & /api/expenses     & Get all expenses \\
POST   & /api/expenses     & Create expense entry \\
PUT    & /api/expenses/:id & Update expense entry \\
DELETE & /api/expenses/:id & Delete expense entry \\
\bottomrule
\end{tabularx}
\caption{Level 1 REST API endpoints for Income and Expense.}
\end{table}

\subsection{Error Handling}
A centralized error-handler middleware returns JSON errors with proper HTTP
status codes (404 for not found, 500 for server errors).

\subsection{Testing and Verification}
The API was tested during development using \textbf{Thunder Client} (a VS Code
extension), and a shareable \textbf{Postman} collection was created for
submission covering every endpoint (GET, POST, PUT, DELETE for both resources,
including error cases). The collection is included at
\texttt{postman/finpilot-level1.postman\_collection.json}; the \texttt{baseUrl}
variable is set to \texttt{http://localhost:5000/api} and requests are run in
order (Create $\rightarrow$ Get $\rightarrow$ Update $\rightarrow$ Delete).

\finfig{../screenshots/level-1/create-income.png}{0.85}{Creating an income entry via the REST API.}
\finfig{../screenshots/level-1/get-income.png}{0.85}{Retrieving all income entries.}
\finfig{../screenshots/level-1/update-income.png}{0.85}{Updating an existing income entry.}
\finfig{../screenshots/level-1/delete-income.png}{0.85}{Deleting an income entry.}
\finfig{../screenshots/level-1/create-expense.png}{0.85}{Creating an expense entry via the REST API.}
\finfig{../screenshots/level-1/error-invalid-id.png}{0.85}{Error handling for an invalid resource ID.}

\subsection{Contribution to FinPilot}
This task delivered the first working slice of FinPilot's backend --- the
Income and Expense data model and CRUD endpoints --- which the frontend, database
integration, and later financial modules extend and build upon.

%==============================================================================
\chapter{Level 2 --- Intermediate}
%==============================================================================

\section{Task 1 --- Frontend Framework}

\subsection{Objective}
Rebuild the FinPilot frontend using a modern JavaScript framework with a
component-based, type-safe architecture, centralized state management, and typed
API calls with loading states.

\subsection{Technologies Used}
\begin{itemize}[leftmargin=1.4em]
  \item \textbf{React 19} + \textbf{TypeScript} --- component-based, type-safe UI.
  \item \textbf{Vite} --- fast dev server and build tool.
  \item \textbf{React Router v7} --- client-side routing and protected routes.
  \item \textbf{Axios} --- HTTP client with request/response interceptors.
\end{itemize}

\subsection{Implementation}
The frontend was built as a Vite + React + TypeScript project under
\texttt{frontend/}. All components use hooks; global authentication state is held
in \texttt{AuthContext} via the Context API. Typed Axios services live in
\texttt{src/api/}, and every form and list shows spinners while requests are
pending. A small set of reusable UI components was created.

\begin{table}[H]
\centering
\begin{tabularx}{\linewidth}{>{\ttfamily\bfseries}l X}
\toprule
\normalfont\bfseries Component & Purpose \\
\midrule
Button  & Variants (primary/secondary/danger/ghost), built-in loading spinner, full-width option \\
Input   & Labeled input with consistent focus styling \\
Card    & Surface container used for panels and auth forms \\
Spinner & Inline loading indicator for lists \\
\bottomrule
\end{tabularx}
\caption{Reusable UI components in \texttt{src/components/ui/}.}
\end{table}

\subsection{State Management and Feedback}
\texttt{AuthContext} centralizes the authenticated user, token, and the
\texttt{login} / \texttt{register} / \texttt{logout} actions. The Dashboard lifts
income and expense totals into local state to render a live
\textbf{Total Income / Total Expenses / Balance} summary. Login and Register
buttons show a spinner and disable while a request is in flight; income and
expense lists show a loading spinner on first fetch, an empty state when there is
no data, and inline error messages on failure.

\subsection{Styling}
A custom design system was added via CSS variables in
\texttt{frontend/src/index.css} with component styles in
\texttt{frontend/src/App.css}, including automatic light/dark mode support.

\subsection{Frontend Project Structure}
\begin{lstlisting}[language=bash,caption={Frontend source layout.}]
frontend/src/
  api/            # Axios client + typed service calls (auth, income, expense)
  components/
    ui/           # Reusable UI: Button, Input, Card, Spinner
    incomeList.tsx
    expenseList.tsx
    ProtectedRoute.tsx
  context/
    AuthContext.tsx   # Global auth state (login, register, logout)
  pages/
    Login.tsx
    Register.tsx
  types/          # Shared TypeScript interfaces
  App.tsx         # Routes + Dashboard
\end{lstlisting}

\subsection{Screenshots}
\finfig{../screenshots/level-2/login.png}{0.8}{Login page built with React and TypeScript.}
\finfig{../screenshots/level-2/register.png}{0.8}{Registration page with client-side validation.}
\finfig{../screenshots/level-2/dashboard.png}{0.85}{Dashboard showing the live income/expense/balance summary.}

\subsection{Contribution to FinPilot}
This task delivered the application's user interface foundation --- routing,
global auth state, typed API clients, and the reusable component library --- that
all later feature pages (budgets, savings, investments, analytics, admin) reuse.

\section{Task 2 --- Authentication}

\subsection{Objective}
Implement user authentication (signup/login) with secure password storage and
stateless sessions, protect backend routes, and restrict access based on user
roles (\texttt{user} / \texttt{admin}).

\subsection{Technologies Used}
\texttt{bcryptjs} for password hashing and \texttt{jsonwebtoken} (JWT) for
stateless sessions, integrated on the frontend through an Axios interceptor and
\texttt{AuthContext}.

\subsection{Auth Endpoints}
\begin{table}[H]
\centering
\begin{tabularx}{\linewidth}{l l l X}
\toprule
\textbf{Method} & \textbf{Endpoint} & \textbf{Access} & \textbf{Description} \\
\midrule
POST & /api/auth/register & Public        & Create account, returns JWT \\
POST & /api/auth/login    & Public        & Login, returns JWT \\
GET  & /api/auth/me       & Authenticated & Get current user \\
GET  & /api/auth/users    & Admin only    & List all users (role-based) \\
\bottomrule
\end{tabularx}
\caption{Authentication endpoints and required access level.}
\end{table}

\subsection{Password Hashing}
Passwords are never stored in plain text. The \texttt{User} model hashes the
password before saving and exposes a \texttt{comparePassword()} helper for login.

\begin{lstlisting}[language=Java,caption={Password hashing in a pre-save hook on the User model.}]
userSchema.pre('save', async function () {
  if (!this.isModified('password')) return;
  const salt = await bcrypt.genSalt(10);
  this.password = await bcrypt.hash(this.password, salt);
});
\end{lstlisting}

\subsection{JWT Flow}
\begin{enumerate}[leftmargin=1.6em]
  \item On register/login the server signs a JWT containing the user \texttt{id}
        and \texttt{role} (expires in 7 days).
  \item The frontend stores the token and attaches it as
        \texttt{Authorization: Bearer <token>} on every request via an Axios
        interceptor.
  \item The \texttt{protect} middleware verifies the token, loads the user, and
        attaches it to \texttt{req.user}.
  \item On a \texttt{401} response the frontend clears the token and redirects to
        \texttt{/login}.
\end{enumerate}

\subsection{Role-Based Access Control}
The \texttt{User} model has a \texttt{role} field
(\texttt{enum: ['user', 'admin']}, default \texttt{user}). The
\texttt{authorize(...roles)} middleware returns \texttt{403 Forbidden} when the
user's role is not permitted --- for example, \texttt{GET /api/auth/users} is
guarded by \texttt{protect, authorize('admin')}. On the frontend,
\texttt{AuthContext} verifies the session on load via \texttt{GET /api/auth/me}
and \texttt{ProtectedRoute} redirects unauthenticated users to \texttt{/login}.

\subsection{Testing and Verification}
A Postman collection at
\texttt{postman/finpilot-level2.postman\_collection.json} covers register, login,
\texttt{me}, protected access, and the admin-only route. The flow verifies that a
normal user token receives \texttt{403} on the admin route while an admin token
receives the user list. Error responses are: \texttt{400} (email already
registered / invalid input), \texttt{401} (missing, invalid, or expired token;
wrong credentials), \texttt{403} (authenticated but insufficient role), and
\texttt{404} (user not found).

\finfig{../screenshots/level-2/auth-register.png}{0.85}{Successful registration returning a JWT.}
\finfig{../screenshots/level-2/auth-login.png}{0.85}{Successful login returning a JWT.}
\finfig{../screenshots/level-2/auth-401.png}{0.85}{Protected route rejecting a request with no token (401).}

\subsection{Contribution to FinPilot}
Authentication is the security backbone of FinPilot: it guarantees that every
private financial record is tied to and accessible only by its owner, and it
gates the admin tools behind role checks.

\section{Task 3 --- Database Integration}

\subsection{Objective}
Integrate MongoDB using the Mongoose ODM, define models and relationships,
enforce data validation before saving, and add indexes for query optimization,
with all CRUD operations scoped to the authenticated user.

\subsection{Technologies Used}
MongoDB (Atlas) with the Mongoose ODM. The connection is defined in
\texttt{src/config/db.js} and schemas live in \texttt{src/models/}.

\subsection{Models}
\begin{table}[H]
\centering
\begin{tabularx}{\linewidth}{>{\bfseries}l X l}
\toprule
Model & Key Fields & Notes \\
\midrule
User    & name, email (unique, lowercase), password (min 6), role (enum user/admin) & Password hashed pre-save \\
Income  & user (ref User), title, amount (min 0), date, category, paymentMethod, description & Indexed by \texttt{\{ user, date \}} \\
Expense & user (ref User), title, amount (min 0), category, date, paymentMethod, description & Indexed by \texttt{\{ user, date \}} \\
\bottomrule
\end{tabularx}
\caption{Core Mongoose models and their key fields.}
\end{table}

\subsection{Relationships}
\texttt{Income} and \texttt{Expense} each hold a \texttt{user} field that
references the \texttt{User} collection via a Mongoose \texttt{ObjectId} ref.
Every controller query filters by \texttt{req.user.id}, so a user can only read or
modify their own records.

\begin{lstlisting}[language=Java,caption={User reference on the Income/Expense schema.}]
user: { type: mongoose.Schema.Types.ObjectId, ref: 'User', required: true }
\end{lstlisting}

\subsection{Data Validation}
Validation runs automatically before documents are saved:
\begin{itemize}[leftmargin=1.4em]
  \item \textbf{Required fields:} \texttt{user}, \texttt{title}, \texttt{amount}.
  \item \textbf{amount:} \texttt{Number} with \texttt{min: 0} --- negative values
        are rejected.
  \item \textbf{String fields:} \texttt{trim: true} removes accidental whitespace.
  \item \textbf{email:} \texttt{unique} + \texttt{lowercase} to prevent duplicate
        accounts.
  \item \textbf{role:} restricted to \texttt{enum: ['user', 'admin']}.
  \item \textbf{password:} \texttt{minlength: 6}.
\end{itemize}

\subsection{Indexing and Optimization}
A compound index \texttt{\{ user: 1, date: -1 \}} on both Income and Expense makes
the common "list my records, newest first" query fast. The \texttt{unique}
constraint on \texttt{email} creates an index that enforces one account per email
and speeds up login lookups.

\subsection{CRUD Endpoints}
\begin{table}[H]
\centering
\begin{tabularx}{\linewidth}{l l X}
\toprule
\textbf{Method} & \textbf{Endpoint} & \textbf{Description} \\
\midrule
GET    & /api/income       & List current user income \\
POST   & /api/income       & Create income entry \\
PUT    & /api/income/:id   & Update income entry \\
DELETE & /api/income/:id   & Delete income entry \\
GET    & /api/expenses     & List current user expenses \\
POST   & /api/expenses     & Create expense entry \\
PUT    & /api/expenses/:id & Update expense entry \\
DELETE & /api/expenses/:id & Delete expense entry \\
\bottomrule
\end{tabularx}
\caption{User-scoped CRUD endpoints backed by MongoDB.}
\end{table}

\subsection{Connection}
\begin{lstlisting}[language=Java,caption={MongoDB connection via Mongoose.}]
mongoose.connect(process.env.MONGO_URI)
\end{lstlisting}
The URI is stored in \texttt{backend/.env} and is never committed.

\subsection{Testing and Verification}
The centralized error handler returns JSON with proper status codes:
\texttt{400} for Mongoose validation errors (missing/invalid fields, negative
amount), \texttt{403} when attempting to modify a record owned by another user,
and \texttt{500} for unexpected server/database errors.

\finfig{../screenshots/level-2/db-connected.png}{0.85}{MongoDB Atlas connected to the application.}
\finfig{../screenshots/level-2/db-create.png}{0.85}{Creating a record that passes schema validation.}
\finfig{../screenshots/level-2/db-validation-error.png}{0.85}{A negative amount rejected by validation.}

\subsection{Contribution to FinPilot}
Database integration turns FinPilot's in-memory operations into durable,
per-user persisted data with validation and fast queries --- the storage
foundation every financial module depends on.

%==============================================================================
\chapter{Level 3 --- Advanced}
%==============================================================================

\section{Task 1 --- Full-Stack Application}

\subsection{Objective}
Extend FinPilot into a complete full-stack (MERN) personal finance application
covering all core modules with full role-based access control (user/admin).

\subsection{Modules Implemented}
\begin{table}[H]
\centering
\begin{tabularx}{\linewidth}{>{\bfseries}l X}
\toprule
Module & Description \\
\midrule
Authentication    & Register, login, JWT, bcrypt password hashing, protected routes \\
Dashboard         & Aggregated stats, charts (income vs expense, category breakdown, budget usage) \\
Income \& Expense & Full CRUD, scoped per user \\
Budgets           & Create/edit/delete, auto-calculated spent amount, exceed alerts \\
Savings Goals     & Create goals, deposit money, auto-complete on target reached \\
Investments       & Track stocks/crypto/gold/etc., profit/loss calculation \\
Transactions      & Unified activity log, search/filter/sort, CSV \& PDF export \\
Categories        & Custom category CRUD per user \\
Analytics         & Charts for income growth, savings growth, category spending \\
Profile           & Update personal info, change password, theme preference \\
Admin Panel       & User management, suspend users, view platform-wide stats \\
\bottomrule
\end{tabularx}
\caption{Full-stack modules implemented in Level 3.}
\end{table}

\subsection{Role-Based Access}
Regular users can only access and modify their own data, enforced via
\texttt{user} ownership checks on every query. Admin-only routes are protected by
the \texttt{authorize('admin')} middleware, covering user management and
platform statistics. (The underlying JWT/bcrypt mechanism is described in
Chapter~3, Task~2.)

\subsection{UI/UX Design}
The application uses a dashboard-style layout with a collapsible, grouped sidebar
navigation instead of a flat link list, to keep the interface manageable as the
feature count grew.
\begin{itemize}[leftmargin=1.4em]
  \item \textbf{Grouped sections:} Overview, Money (Income/Expenses/Budgets),
        Growth (Savings/Investments), Records (Transactions/Categories), Insights
        (Analytics), and a pinned Account section (Profile/Admin/Logout).
  \item Sidebar collapses to an icon-only rail with flyout popovers, expandable on
        demand; the active route auto-expands its parent group and is highlighted.
  \item Fully responsive: overlay drawer navigation on mobile.
  \item Reusable Modal component for create/edit forms (Budgets, Investments,
        Categories, Savings deposits) instead of always-visible inline forms.
  \item Toast notification system for form feedback and real-time Socket.io
        alerts.
  \item Skeleton loading states and empty-state messaging on all list views.
  \item Light/dark theme support via CSS custom properties, respected across every
        component.
  \item Consistent Card and Badge components for stats, budget progress, and
        status indicators.
\end{itemize}

\subsection{Screenshots}
\finfig{../screenshots/level-3/dashboard.png}{0.9}{Dashboard overview with charts and aggregated stats.}
\finfig{../screenshots/level-3/sidebar-collapsed.png}{0.9}{Collapsed icon-only sidebar navigation.}
\finfig{../screenshots/level-3/budgets.png}{0.9}{Budgets with progress bars and exceed alerts.}
\finfig{../screenshots/level-3/savings.png}{0.9}{Savings goals with goal tracking and deposits.}
\finfig{../screenshots/level-3/investments.png}{0.9}{Investments portfolio with profit/loss.}
\finfig{../screenshots/level-3/transactions.png}{0.9}{Transactions with search, filter, and export.}
\finfig{../screenshots/level-3/admin-panel.png}{0.9}{Admin panel with user management and platform stats.}
\finfig{../screenshots/level-3/dark-theme.png}{0.9}{Light/dark mode support across the application.}

\subsection{Deployment Readiness}
The backend is configured for Render (env vars: \texttt{MONGO\_URI},
\texttt{JWT\_SECRET}, \texttt{FRONTEND\_URL}, \texttt{PORT}); the frontend is
configured for Vercel with SPA rewrite rules; and the database uses MongoDB Atlas
(cloud-hosted, in use since Level 1).

\subsection{Contribution to FinPilot}
This task integrates every prior building block --- auth, UI framework, and
database --- into the complete FinPilot product, adding the budgets, savings,
investments, transactions, analytics, and admin modules that define the app.

\section{Task 2 --- WebSockets}

\subsection{Objective}
Deliver private, real-time notifications to users using WebSockets so that
finance events (such as budget alerts and savings-goal completion) reach the user
instantly without a page refresh.

\subsection{Why WebSockets Are Used}
FinPilot generates events that should reach the user immediately --- for example,
a budget being exceeded or a savings goal being completed. Polling the REST API
for these would be wasteful and slow, so a persistent Socket.io channel is used to
push notifications to the specific user the moment they occur.

\subsection{Technologies Used}
Socket.io on the backend (attached to the same HTTP server as the REST API) and
the Socket.io client on the React frontend.

\subsection{How the Client Communicates with the Backend}
Socket.io is attached to the HTTP server at startup, with CORS restricted to the
frontend origin. Each client, after authenticating, emits a \texttt{join} event
with its user ID to join a private room named after that user ID. The backend then
emits notification events only to that room, keeping every user's notifications
private.

\begin{lstlisting}[language=Java,caption={Backend: a per-user private room on connection (socket.js).}]
io.on('connection', (socket) => {
  // Client emits `join` with their userId; join a private room for that user.
  socket.on('join', (userId) => {
    if (userId) socket.join(String(userId));
  });
});

// Safe emit helper: no-ops if io is not initialized (e.g. during tests).
function emitToUser(userId, event, payload) {
  if (!io) return;
  io.to(String(userId)).emit(event, payload);
}
\end{lstlisting}

\subsection{How Real-Time Notifications Work}
When the backend needs to notify a user, it persists a notification and pushes it
to that user's room in a single helper. The client listens for the
\texttt{notification} event and updates the notification bell and toast alerts in
real time.

\begin{lstlisting}[language=Java,caption={Backend: persist and push a notification (notify.js).}]
async function notify(userId, message, type) {
  const notification = await Notification.create({ user: userId, message, type });
  emitToUser(userId, 'notification', notification);
  return notification;
}
\end{lstlisting}

\begin{lstlisting}[language=Java,caption={Frontend: join the room and react to notification events (SocketContext).}]
const socket = io(SOCKET_URL, { transports: ['websocket', 'polling'] });
socket.emit('join', user.id);
socket.on('notification', (payload) => {
  setNotifications((prev) => [payload, ...prev]);
  showToast(payload.message, /* type derived from payload.type */);
});
// ...
return () => { socket.disconnect(); };
\end{lstlisting}

On the client, the toast severity is derived from the notification type
(\texttt{budget-alert} $\rightarrow$ warning, \texttt{goal-completed}
$\rightarrow$ success, otherwise info). The socket connects only when the user is
authenticated and disconnects on logout.

\subsection{Testing and Verification}
The implementation was verified to confirm that:
\begin{itemize}[leftmargin=1.4em]
  \item Notifications are delivered instantly without a page refresh.
  \item Notifications are private --- only the relevant user's room receives their
        own events.
  \item The socket disconnects cleanly on logout.
\end{itemize}

\finfig{../screenshots/level-3/live-notification.png}{0.85}{Live notification toast delivered in real time.}
\finfig{../screenshots/level-3/notification-bell.png}{0.85}{Notification bell dropdown listing user notifications.}
\finfig{../screenshots/level-3/socket-connection.png}{0.85}{Active Socket.io connection shown in browser devtools.}

\subsection{Contribution to FinPilot}
WebSockets make FinPilot feel live: budget alerts, savings-goal completion, and
other events surface to the correct user immediately, turning the finance data
into an interactive, responsive experience.

%==============================================================================
\chapter{API Reference}
%==============================================================================
The complete set of REST endpoints exposed by the FinPilot backend is summarized
below. Endpoints marked \emph{Required} need a valid JWT; \emph{Admin} endpoints
additionally require the admin role; \emph{Public} endpoints are accessible
without authentication.

\begin{longtable}{l L{4.9cm} L{5.2cm} l}
\toprule
\textbf{Method} & \textbf{Endpoint} & \textbf{Description} & \textbf{Auth} \\
\midrule
\endfirsthead
\toprule
\textbf{Method} & \textbf{Endpoint} & \textbf{Description} & \textbf{Auth} \\
\midrule
\endhead
\bottomrule
\endfoot
POST & /api/auth/register & Create account, returns token and user & Public \\
POST & /api/auth/login & Login, returns token and user & Public \\
GET & /api/auth/me & Get current authenticated user & Required \\
GET & /api/auth/users & List users (admin only) & Admin \\
GET/POST & /api/income & List / create incomes (scoped to user) & Required \\
PUT/DELETE & /api/income/:id & Update / delete an income owned by the user & Required \\
GET/POST & /api/expenses & List / create expenses (scoped to user) & Required \\
PUT/DELETE & /api/expenses/:id & Update / delete an expense owned by the user & Required \\
GET/POST & /api/budgets & List / create budgets & Required \\
PUT/DELETE & /api/budgets/:id & Update / delete budget & Required \\
GET/POST & /api/savings & List / create savings goals & Required \\
PUT/DELETE & /api/savings/:id & Update / delete savings goal & Required \\
POST & /api/savings/:id/deposit & Deposit to a savings goal & Required \\
GET/POST & /api/investments & List / create investments & Required \\
PUT/DELETE & /api/investments/:id & Update / delete investment & Required \\
GET/POST & /api/categories & List / create categories & Required \\
PUT/DELETE & /api/categories/:id & Update / delete category & Required \\
GET & /api/transactions & List transactions (filters supported) & Required \\
GET & /api/transactions/export/csv & Export transactions as CSV & Required \\
GET & /api/transactions/export/pdf & Export transactions as PDF & Required \\
GET & /api/notifications & List notifications for user & Required \\
DELETE & /api/notifications/:id & Delete a notification & Required \\
PUT & /api/notifications/:id/read & Mark notification read & Required \\
GET & /api/dashboard/stats & Dashboard stats (guest-safe payload if unauthenticated) & Public \\
GET/PUT & /api/profile & Get / update profile & Required \\
PUT & /api/profile/change-password & Change user password & Required \\
GET/DELETE/PUT & /api/admin/* & Admin reports, list/delete users, suspend users & Admin \\
\caption{FinPilot REST API reference.}
\end{longtable}

%==============================================================================
\chapter{Security}
%==============================================================================
FinPilot implements the following security mechanisms, each drawn from the actual
implementation:
\begin{itemize}[leftmargin=1.4em]
  \item \textbf{JWT authentication.} Stateless sessions using signed JSON Web
        Tokens that carry the user \texttt{id} and \texttt{role} and expire in 7
        days.
  \item \textbf{bcrypt password hashing.} Passwords are hashed with
        \texttt{bcryptjs} in a pre-save hook and never stored in plain text; a
        \texttt{comparePassword()} helper validates logins.
  \item \textbf{Protected backend routes.} The \texttt{protect} middleware
        validates the Bearer token on every private route and populates
        \texttt{req.user}.
  \item \textbf{Role-based access control.} The \texttt{authorize(...roles)}
        middleware guards admin-only endpoints, returning \texttt{403 Forbidden}
        for insufficient roles.
  \item \textbf{User-specific data access.} Every controller query filters by
        \texttt{req.user.id}, so users can only read or modify their own records.
  \item \textbf{CORS restriction.} Both the HTTP API and the Socket.io channel
        restrict origins to the configured frontend URL.
\end{itemize}

%==============================================================================
\chapter{Deployment}
%==============================================================================
FinPilot is deployed as a three-part cloud application:
\begin{table}[H]
\centering
\begin{tabularx}{\linewidth}{>{\bfseries}l X}
\toprule
Component & Platform \\
\midrule
Frontend & Vercel (SPA rewrite rules configured) \\
Backend  & Render (\texttt{render.yaml} blueprint; env vars \texttt{MONGO\_URI}, \texttt{JWT\_SECRET}, \texttt{FRONTEND\_URL}, \texttt{PORT}) \\
Database & MongoDB Atlas (cloud-hosted, in use since Level 1) \\
\bottomrule
\end{tabularx}
\caption{Deployment targets for each application layer.}
\end{table}

\noindent\textbf{Live URLs}
\begin{itemize}[leftmargin=1.4em]
  \item Frontend: \url{https://finpilot-oiek-q4zxkmfqt-sushil-bhattas-projects.vercel.app/}
  \item Backend API: \url{https://finplot-rhir.onrender.com}
\end{itemize}

%==============================================================================
\chapter{Final Project Summary}
%==============================================================================

\section{What Was Built}
FinPilot is a complete full-stack personal finance and investment dashboard.
Starting from a verified development environment and a simple Income/Expense REST
API, the project grew into a React + TypeScript frontend, a JWT/bcrypt
authentication system with role-based access control, a validated and indexed
MongoDB data layer, and a full suite of finance modules --- dashboard, income,
expenses, budgets, savings goals, investments, transactions (with CSV/PDF
export), categories, analytics, profile, and an admin panel --- topped with
real-time Socket.io notifications.

\section{Technologies Learned and Applied}
React 19, TypeScript, Vite, React Router, Axios, and Recharts on the frontend;
Node.js, Express, JWT, bcryptjs, Socket.io, and PDFKit on the backend; and MongoDB
with Mongoose for persistence --- deployed across Vercel, Render, and MongoDB
Atlas.

\section{Full-Stack Concepts Demonstrated}
\begin{itemize}[leftmargin=1.4em]
  \item RESTful API design with full CRUD and centralized error handling.
  \item Component-based, type-safe frontend architecture with global state and
        typed API clients.
  \item Secure authentication, password hashing, and role-based authorization.
  \item Relational-style data modeling in MongoDB with validation, references,
        and indexing, scoped per user.
  \item Real-time client--server communication over WebSockets with private
        per-user rooms.
  \item Cloud deployment of a three-tier application.
\end{itemize}

\section{Deployment Status}
The application is deployed and accessible: the frontend on Vercel, the backend
on Render, and the database on MongoDB Atlas.

\section{Overall Outcome}
All seven internship tasks across Levels 1--3 were completed, resulting in a
cohesive, deployed full-stack application. FinPilot demonstrates end-to-end
ownership of a modern web product --- from environment setup and API design to
authentication, data modeling, a polished UI, real-time features, and cloud
deployment.

\end{document}
